\section{Proofs}\label{app:proofs}

\paragraph{Proof of \cref{prop:optimal-level}.}
Conditional on $\F_t$, $P_{t+1}=P_te^{r_{t+1}}$ and $P_t$ is known. For any $\F_t$-measurable $f$,
$\E[(P_{t+1}-f)^2\mid\F_t]=\Var(P_{t+1}\mid\F_t)+(\E[P_{t+1}\mid\F_t]-f)^2$, minimised at
$f^\ast=\E[P_{t+1}\mid\F_t]=P_tM_t(1)$ with value
$\Var(P_{t+1}\mid\F_t)=P_t^2\{M_t(2)-M_t(1)^2\}$. For $f=P_t$ the value is
$P_t^2\E[(e^{r}-1)^2\mid\F_t]=P_t^2\{M_t(2)-2M_t(1)+1\}$; subtracting gives \eqref{eq:mse-gap}.\qed

\paragraph{Proof of \cref{prop:bound}.}
Divide \eqref{eq:mse-gap} by the random-walk MSE: the ratio is
$\rho_t=(M_1-1)^2/(M_2-2M_1+1)$. Under normality $M_1=e^{m}$ and $M_2=e^{2m+\sigma^2}$ with
$m=\mu+\sigma^2/2$. Expanding, $M_1-1=m+O(m^2)$ and
$M_2-2M_1+1=e^{2m}(e^{\sigma^2}-1)+(e^{m}-1)^2=\sigma^2+m^2+O(\sigma^3)$ since $m=O(\sigma)$ for
bounded $s$. Hence $\rho_t=m^2/(\sigma^2+m^2)+O(\sigma^3)=s_t^2/(1+s_t^2)+O(\sigma^3)$. For the
unconditional statement write the optimal and RW MSEs as $\E[P_t^2\{\cdot\}]$; the ratio of
expectations is the $w_t$-weighted average of $\rho_t$ with $w_t\propto P_t^2(\sigma_t^2+m_t^2)$,
and $\mathrm{RMSE}$ ratios are square roots of MSE ratios. The approximation uses
$\sqrt{1-x}\approx1-x/2$.\qed

\paragraph{Proof of \cref{prop:da}.}
If $f_{t+1}=P_t$ then $f_{t+1}-f_t=P_t-P_{t-1}$, whose sign equals that of $r_t$, while
$P_{t+1}-P_t$ has the sign of $r_{t+1}$. The expectation follows from
$\1\{\operatorname{sign}a=\operatorname{sign}b\}=\tfrac12(1+\operatorname{sign}a\operatorname{sign}b)$ for non-zero $a,b$. If
$f_{t+1}=P_t+\varepsilon_{t+1}$ then $f_{t+1}-f_t=(P_t-P_{t-1})+(\varepsilon_{t+1}-\varepsilon_t)$,
which has the sign of $r_t$ whenever $|\varepsilon_{t+1}-\varepsilon_t|<|P_t-P_{t-1}|$; the stated
condition is the same inequality in log units to first order.\qed

\paragraph{Proof of \cref{prop:ident}.}
Fix covariates $(\bx,\bv,\bz)$ in the open set of \cref{ass:ident}(a). Finite mixtures of
univariate normals are identifiable \citep{teicher1963identifiability}: equality of the mixtures
implies equality of the sets $\{(\pi_k,\mu_k,\sigma_k^2)\}_{k}$ and $\{(\pi'_k,\mu'_k,\sigma'^2_k)\}_k$,
which by \cref{ass:ident}(b) have $K$ distinct elements on a set of positive probability. Hence there
is a permutation $\varsigma_{(\bx,\bv,\bz)}$ with
$(\pi'_{\varsigma(k)},\mu'_{\varsigma(k)},\sigma'^2_{\varsigma(k)})=(\pi_k,\mu_k,\sigma_k^2)$. All
component parameters are continuous in the covariates and distinct, so $\varsigma$ is locally
constant and therefore constant on the connected open set. Suppose $\varsigma(k)=l\neq k$. Then
$\bm\beta_k^\top\tilde\bx^{(k)}=\bm\beta_l'^\top\tilde\bx^{(l)}$ on an open set. By
\cref{ass:ident}(c) the two covariate sets differ, so one side varies with a covariate on which the
other does not depend; non-singularity of the designs forces the corresponding slope to be zero,
contradicting (c), unless both experts have no slopes, which (c) excludes for more than one expert.
Hence $\varsigma$ is the identity. Equality of $\mu_k$, $\log\sigma_k^2$ and of the log-odds
$\log(\pi_k/\pi_K)=\bm\gamma_k^\top\tilde\bz$ as linear functions on an open set, together with
non-singular designs, gives $\bm\beta_k=\bm\beta'_k$, $\bm\delta_k=\bm\delta'_k$ and
$\bm\gamma_k=\bm\gamma'_k$.\qed

\paragraph{Proof of \cref{prop:ergodic}.}
Write $Y_t=(y_t,\dots,y_{t-p+1})$. The chain is a $\varphi$-irreducible, aperiodic $T$-chain because
the transition density is a positive Gaussian mixture, and compact sets are small
\citep[Ch.~5--6]{meyn2009markov}. Let $\rho=\sup_{\bz}\sum_k\pi_k(\bz)\sum_j|\phi_{kj}|<1$ and choose
$V(Y)=1+\sum_{j=1}^pc_j|y_{t-j+1}|$ with $c_1=1$ and $c_j$ decreasing geometrically so that
$c_{j+1}+\rho_j\le c_j\rho'$ for some $\rho<\rho'<1$, where $\rho_j=\sup_{\bz}\sum_k\pi_k|\phi_{kj}|$.
Then
$\E[|y_{t+1}|\mid Y_t]\le\sum_k\pi_k(|\phi_{k0}|+\sum_j|\phi_{kj}||y_{t-j+1}|)+\bar\sigma\sqrt{2/\pi}$,
so $\E[V(Y_{t+1})\mid Y_t]\le\rho'V(Y_t)+C$ for a constant $C$. The Foster--Lyapunov criterion
\citep[Thm.~15.0.1]{meyn2009markov} yields geometric ergodicity with
$\E_{\text{stat}}V<\infty$.\qed

\paragraph{Proof of \cref{prop:closed}.}
Mean, cdf and probability statements are linear in the mixture weights. The variance follows from
the law of total variance with the component label as conditioning variable. For ES, write
$\E[y\1\{y\le q\}]=\sum_k\pi_k\E[y\1\{y\le q\}\mid k]$ and use
$\E[X\1\{X\le q\}]=\mu\Phi(z)-\sigma\varphi(z)$ for $X\sim\N(\mu,\sigma^2)$, $z=(q-\mu)/\sigma$.
For CRPS use $\mathrm{CRPS}(F,y)=\E|X-y|-\tfrac12\E|X-X'|$ with $X,X'\stackrel{iid}{\sim}F$, and
$\E|Z|=A(m,s^2)$ for $Z\sim\N(m,s^2)$ \citep{grimit2006continuous}. Uniqueness of $q_\alpha$ holds
because the mixture cdf is continuous and strictly increasing.\qed

\paragraph{Proof of \cref{prop:monotone}.}
For any $\btheta$, Jensen's inequality gives
$\ell^\lambda_n(\btheta)-\ell^\lambda_n(\btheta^{(s)})\ge Q(\btheta\mid\btheta^{(s)})-Q(\btheta^{(s)}\mid\btheta^{(s)})$
\citep{dempster1977maximum}. Each CM step does not decrease $Q$: CM-step 1 maximises a strictly
concave quadratic exactly; CM-steps 2 and 3 take Newton steps on strictly concave objectives with step
halving that accepts only non-decreasing values. Hence $Q(\btheta^{(s+1)}\mid\btheta^{(s)})\ge
Q(\btheta^{(s)}\mid\btheta^{(s)})$ and the claim follows; the final L-BFGS step is accepted only if it
increases $\ell^\lambda_n$. Convergence of limit points to stationary points follows from
\citet[Thm.~1]{wu1983convergence} applied to the generalised EM map.\qed

\paragraph{Proof of \cref{prop:elasticity}.}
For the softmax with $\bm\gamma_K=\bm0$,
$\partial\pi_k/\partial z_j=\pi_k(\gamma_{kj}-\sum_l\pi_l\gamma_{lj})$. Since $z_j$ enters neither
$\mu_k$ nor $\sigma_k$, $\partial\bar\mu/\partial z_j=\sum_k\mu_k\partial\pi_k/\partial z_j$ and
$\partial P(y>0)/\partial z_j=\sum_k\Phi(\mu_k/\sigma_k)\partial\pi_k/\partial z_j$.\qed

\paragraph{Proof of \cref{prop:fidelity}.}
$\E\log p(y_t\mid\cdot)=\E\int p_0\log p=\E\int p_0\log p_0-\E\int p_0\log(p_0/p)=-H(p_0)-\E\KL(p_0\|p)$.
The entropy cancels in differences.\qed

\paragraph{Proof of \cref{thm:aci}.}
Summing \eqref{eq:aci} gives $\alpha_{T+1}=\alpha_1+\eta\sum_{t=1}^T(a-\mathrm{err}_t)$. If
$\alpha_t\le0$ then $q_t=-\infty$, $\mathrm{err}_t=0$ and $\alpha_{t+1}=\alpha_t+\eta a>\alpha_t$; if
$\alpha_t\ge1$ then $\mathrm{err}_t=1$ and $\alpha_{t+1}<\alpha_t$. Each step moves $\alpha$ by at
most $\eta$, so $\alpha_t\in[-\eta,1+\eta]$ for all $t$ when $\alpha_1\in(0,1)$. Therefore
$|\sum_t(a-\mathrm{err}_t)|=|\alpha_{T+1}-\alpha_1|/\eta\le(\max(\alpha_1,1-\alpha_1)+\eta)/\eta$,
and dividing by $T$ gives the bound \citep[cf.][Prop.~4.1]{gibbs2021adaptive}.\qed

\paragraph{Student-$t$ densities as Gaussian mixtures.}
All benchmarks are scored through the closed forms of \cref{prop:closed}, which requires a Gaussian
mixture representation. For the Student-$t$ benchmarks we use the scale-mixture identity
$t_\nu(\mu,s)=\int_0^\infty\N(\mu,s^2/\lambda)\,g(\lambda)\,d\lambda$ with
$\lambda\sim\mathrm{Gamma}(\nu/2,\text{rate }\nu/2)$. Substituting $u=\nu\lambda/2$,
\begin{equation}
  \E[h(\lambda)]=\frac{1}{\Gamma(\nu/2)}\int_0^\infty h(2u/\nu)\,u^{\nu/2-1}e^{-u}\,du
  \approx\sum_{i=1}^{N}\omega_i\,h(2u_i/\nu),
\end{equation}
where $(u_i,\tilde\omega_i)$ are the nodes and weights of $N$-point generalised Gauss--Laguerre
quadrature with parameter $\alpha=\nu/2-1$ and $\omega_i=\tilde\omega_i/\Gamma(\nu/2)$. The rule is
exact for polynomials of degree $2N-1$ in $u$; for the smooth integrands $h(\lambda)=\varphi(y;\mu,s^2/\lambda)$
the error decays in $N$ at fixed $y$, more slowly for heavier tails. We use $N=64$. The maximal
relative density error on $|y-\mu|\le4s$ is $0.27\%$ for $\nu=2.5$, $0.11\%$ for $\nu=3$, $0.02\%$
for $\nu=4$ and below $0.005\%$ for $\nu\ge5$ (computed against the exact $t$ density in the
replication package), an order of magnitude below the log-score differences we interpret. An
earlier version with $N=16$ had errors of up to $2.9\%$ at $\nu=3$; we report this because the choice
of $N$ matters for heavy-tailed benchmarks.\qed

\paragraph{Proof of \cref{prop:shrink}.}
By \eqref{eq:var}, $V_t=W_t+B_t=W_t(1+B_t/W_t)$, so $w^\ast_t=\bar\mu_t/(\gamma V_t)$ factorises as
stated. $B_t\ge0$ with equality iff $\mu_{kt}=\bar\mu_t$ for all $k$ with $\pi_{kt}>0$; a
mean-preserving spread of $(\mu_{kt})$ increases $B_t$ and leaves $\bar\mu_t$ and $W_t$ unchanged.\qed

\paragraph{Proof of \cref{prop:band}.}
The objective $J(w)=w\bar\mu-\tfrac\gamma2w^2V-\kappa|w-w_-|$ is strictly concave. Its
superdifferential is $\bar\mu-\gamma Vw-\kappa\,\partial|w-w_-|$. If $|w_--w^\ast|\le b$ then $0$
lies in the superdifferential at $w=w_-$ (take the subgradient $(\bar\mu-\gamma Vw_-)/\kappa\in[-1,1]$).
If $w_-<w^\ast-b$ the optimum is at $w>w_-$ where the first-order condition
$\bar\mu-\gamma Vw-\kappa=0$ gives $w=w^\ast-b$; symmetrically for $w_->w^\ast+b$. Concavity makes
projection onto the box optimal for the constrained problem.\qed

\paragraph{Proof of \cref{prop:ident-shared}.}
The first part of the proof of \cref{prop:ident} does not use \cref{ass:ident}(c): it yields a
permutation $\varsigma$, constant on a connected open set of covariates, with
$(\pi'_{\varsigma(k)},\mu'_{\varsigma(k)},\sigma'^2_{\varsigma(k)})=(\pi_k,\mu_k,\sigma_k^2)$. For a pair of
experts with different covariate sets the argument of \cref{prop:ident} still excludes
$\varsigma(k)=l$. Hence $\varsigma$ can only permute experts within a shared group. Suppose it
exchanges $k$ and $l$ with $\delta_{k,0}<\delta_{l,0}$. Equality of the log variances as linear
functions of $\tilde\bv$ on an open set gives $\bm\delta'_l=\bm\delta_k$ and $\bm\delta'_k=\bm\delta_l$,
so $\delta'_{k,0}=\delta_{l,0}>\delta_{k,0}=\delta'_{l,0}$, which violates the convention for
$\btheta'$. Thus $\varsigma$ is the identity and the remainder of the proof of \cref{prop:ident}
applies. The convention requires $\delta_{k,0}\neq\delta_{l,0}$; if the two variance functions
coincided, an ordering of the mean intercepts would serve the same purpose.\qed

\paragraph{Proof of \cref{prop:lm-gate}.}
The score expression \eqref{eq:lm-score} is \eqref{eq:score} for the gate block with $\tilde\bz_t$
replaced by the augmented vector, evaluated at $\bm\psi=\bm0$, where the prior probabilities
$\pi_{tk}$ do not depend on $\bm w_t$. Write $\bm\vartheta=(\btheta,\bm\psi)$,
$\bm s_n(\bm\vartheta)=\sum_t\partial\ell_t(\bm\vartheta)$ and $A$ for the information matrix (minus the
expected Hessian) of the augmented model at $\bm\vartheta^\ast=(\btheta^\ast,\bm0)$. Under $H_0$,
$\bm\vartheta^\ast$ is the pseudo-true value, so $\E\,\partial\ell_t(\bm\vartheta^\ast)=\bm0$ and, by the
central limit theorem for $\beta$-mixing sequences under \cref{ass:asym},
$n^{-1/2}\bm s_n(\bm\vartheta^\ast)\xrightarrow{d}\N(\bm0,B)$. The restricted estimator satisfies
$\sqrt n(\hat\btheta-\btheta^\ast)=A_{\theta\theta}^{-1}n^{-1/2}\bm s_{n,\theta}(\bm\vartheta^\ast)+o_p(1)$
by \cref{thm:qmle}; the penalty contributes $O_p(n^{-1/2})$. A mean-value expansion gives
\[
  n^{-1/2}\bm s_n(\hat\btheta,\bm0)=n^{-1/2}\bm s_n(\bm\vartheta^\ast)-
  \begin{pmatrix}A_{\theta\theta}\\A_{\psi\theta}\end{pmatrix}\sqrt n(\hat\btheta-\btheta^\ast)+o_p(1).
\]
Pre-multiplying by $C=[-A_{\psi\theta}A_{\theta\theta}^{-1},\,I]$ annihilates the second term, because
$-A_{\psi\theta}A^{-1}_{\theta\theta}A_{\theta\theta}+A_{\psi\theta}=0$. Hence
$C\,n^{-1/2}\bm s_n(\hat\btheta,\bm0)=C\,n^{-1/2}\bm s_n(\bm\vartheta^\ast)+o_p(1)\xrightarrow{d}\N(\bm0,CBC^\top)$.
$\hat A$ and the Newey--West $\hat B$ are consistent under the moment and mixing conditions of
\cref{ass:asym} \citep{newey1987simple}, $C\hat BC^\top$ is non-singular because $B$ is, and the
continuous mapping theorem gives the $\chi^2_{(K-1)d}$ limit. Under $\bm\psi^\ast=\bm h/\sqrt n$ the
same expansion adds $A_{\cdot\psi}\bm h$ to the limit, and $CA_{\cdot\psi}=A_{\psi\psi\cdot\theta}\coloneqq
A_{\psi\psi}-A_{\psi\theta}A_{\theta\theta}^{-1}A_{\theta\psi}$, so the non-centrality is
$\bm h^\top A_{\psi\psi\cdot\theta}(CBC^\top)^{-1}A_{\psi\psi\cdot\theta}\bm h$. For the Wald statistic,
the $\bm\psi$ rows of $A^{-1}$ equal $A_{\psi\psi\cdot\theta}^{-1}C$ by the partitioned-inverse formula,
so the $\bm\psi$ block of the sandwich $A^{-1}BA^{-1}$ is
$A_{\psi\psi\cdot\theta}^{-1}CBC^\top A_{\psi\psi\cdot\theta}^{-1}$ and $\sqrt n\hat{\bm\psi}$ has
mean $\bm h$; the two non-centrality parameters coincide.\qed

\paragraph{Proof of \cref{lem:logvar-cost}.}
$\E[-\log\varphi(y;\mu,s^2)]=\tfrac12\log(2\pi s^2)+\sigma^2/(2s^2)$. Subtracting the same expression
at $s^2=\sigma^2$ gives $\tfrac12\{\log(s^2/\sigma^2)+\sigma^2/s^2-1\}=\tfrac12(u+e^{-u}-1)$. The QLIKE
loss $\sigma^2/s^2-\log(\sigma^2/s^2)-1$ equals $e^{-u}+u-1$. Expanding $e^{-u}$ gives the series.
For random $u$, take expectations of $u^2/4$: $\E u^2=b^2+v$.\qed

\paragraph{Proof of \cref{prop:eiv}.}
The slopes of a linear projection are $\Var(\bv_t)^{-1}\Cov(\bv_t,T_t)$ for target $T_t$. Independence
gives $\Var(\bv_t)=\Sigma^\ast+\Sigma_\xi$ and $\Cov(\bv_t,T_t)=\Cov(\bv^\ast_t,T_t)=\Sigma^\ast\bm\delta$,
the last equality being the normal equations of the latent projection. For
$\Sigma_\xi=\omega^2\bm e_1\bm e_1^\top$ the Sherman--Morrison formula gives
$(\Sigma^\ast+\omega^2\bm e_1\bm e_1^\top)^{-1}\Sigma^\ast\bm\delta
=\bm\delta-\omega^2\delta_1\Sigma^{\ast-1}\bm e_1/(1+\omega^2[\Sigma^{\ast-1}]_{11})$.
By the partitioned-inverse formula, $\Sigma^{\ast-1}\bm e_1=\nu_1^{-1}(1,-\bm b^\top)^\top$, so
$[\Sigma^{\ast-1}]_{11}=\nu_1^{-1}$ and
$\bm\delta^{\mathrm{obs}}=\bm\delta-\frac{\omega^2\delta_1}{\nu_1+\omega^2}(1,-b_5,-b_{22})^\top$, which is
\eqref{eq:eiv}. Summing the three components gives the change in the sum of slopes.\qed
